\documentclass[12pt]{article}
\usepackage[utf8]{inputenc}
\usepackage[T1]{fontenc}
\usepackage{hyperref}
\title{Kaiy Language Resources for K-12 Curriculum Development}
\author{Compiled by Claude Sonnet 5 (claude-sonnet-5) for Andr\'as Kornai}
\date{}
\begin{document}
\maketitle
\section{Introduction}
Kaiy (also Kai, Taori-Kei, Taori-Kaiy, Todi) is a Lakes Plain language of Rufaer District, Mamberamo Raya Regency, Papua, Indonesia, spoken in Kaiy and Kokou villages. Estimates are 220 (2000) to 250 speakers. Glottocode kaiy1239.
\section{Primary Educational Resources}
None found.
\section{Linguistic and Reference Materials}
\begin{itemize}
\item Wikipedia, Kaiy language: \url{https://en.wikipedia.org/wiki/Kaiy_language}
\item ASJP word list: \url{https://asjp.clld.org/languages/KAIY}
\item Rosetta Project entry: \url{https://rosettapanglossia.longnow.org/index.php?title=Tcq}
\end{itemize}
\section{Religious and Community Resources}
\begin{itemize}
\item Joshua Project, Taori-Kei, Kaiy: \url{https://joshuaproject.net/people_groups/15261}
\end{itemize}
Joshua Project reports no complete Bible translation.
\section{Digital and Audio Resources}
Global Recordings lists no recordings: \url{https://globalrecordings.net/en/language/tcq}
\section{Curriculum Development Context}
Almost nothing exists beyond word lists and reference entries. A curriculum would require new documentation work.
\end{document}